\paragrafos{Απόλυτη τιμή, αντίθετοι και σύγκριση: ποιος είναι πιο κρύος;}{Α.7.2, «Απόλυτη τιμή ρητού – Αντίθετοι ρητοί – Σύγκριση ρητών»}

\begin{anagke}
Στο πούλμαν της επιστροφής η παρέα συγκρίνει τις θερμοκρασίες του πίνακα του καιρού: Ιωάννινα $-3$, Φλώρινα $-10$. «Ποια πόλη είναι πιο κρύα;» ρωτάει η Ρήνη. Ο Βασίλης απαντάει αμέσως: «Το δέκα είναι μεγαλύτερο από το τρία, άρα το $-10$ είναι μεγαλύτερο από το $-3$. Άρα η Φλώρινα είναι πιο ζεστή». Ο Νίκος διαμαρτύρεται: «Δεν μπορεί να είναι πιο ζεστό το δέκα βαθμούς κάτω από το μηδέν από το τρεις βαθμούς κάτω!». Η Αγγελική παίρνει ένα χαρτί και σχεδιάζει τον άξονα: «Όσο πιο αριστερά, τόσο πιο κρύο. Το $-10$ είναι πολύ πιο αριστερά από το $-3$».

Ο Ηλίας προσθέτει ένα δεύτερο ερώτημα: «Αν η Ζωή μένει $3$ στάσεις αριστερά από το σχολείο και ο Στέλιος $3$ στάσεις δεξιά, ποιος μένει πιο μακριά;». Οι δύο μένουν το ίδιο μακριά! Το «πόσο μακριά» δεν εξαρτάται από την κατεύθυνση, αλλά το «πού ακριβώς» εξαρτάται. Έτσι έχουμε δύο έννοιες: την απόσταση από το $0$, που λέγεται απόλυτη τιμή, και τη σχετική θέση δεξιά ή αριστερά, που μας βοηθάει να συγκρίνουμε.
\end{anagke}

\begin{lexi}{απόλυτη τιμή, αντίθετος}{απόλυτη τιμή, αντίθετος}
\lexiword{Απόλυτη} σημαίνει «λυμένη από τους δεσμούς» (από το «απολύω»): η απόλυτη τιμή είναι ο αριθμός «απελευθερωμένος» από το πρόσημό του, δηλαδή μόνο η απόσταση από το $0$. Οι \lexiword{αντίθετοι} είναι αριθμοί σε «αντίθετη» κατεύθυνση: ίδια απόσταση από το $0$, αλλά η μία δεξιά και η άλλη αριστερά.
\end{lexi}

\begin{metafrasi}
Η απόσταση του $-3$ από το $0$. & $|-3|=3$\\
Ο αντίθετος του $5$. & $-5$\\
Ο αντίθετος του $-2{,}7$. & $-(-2{,}7)=+2{,}7$\\
Το $-10$ είναι πιο αριστερά από το $-3$. & $-10<-3$\\
Δύο αριθμοί με απόλυτη τιμή $2$. & $+2$ και $-2$\\
\end{metafrasi}

\tmima{Η θεωρία}

\begin{orismos}{απόλυτη τιμή}
\textbf{Απόλυτη τιμή} ενός ρητού $\alpha$ είναι η απόσταση του σημείου με τετμημένη $\alpha$ από την αρχή $O$ του άξονα. Συμβολίζεται $|\alpha|$.
\end{orismos}

\begin{itemize}[leftmargin=1.2em,itemsep=0.3ex,topsep=2pt]
\item Η απόλυτη τιμή ενός θετικού είναι ο ίδιος ο αριθμός: $|+9{,}63|=9{,}63$.
\item Η απόλυτη τιμή ενός αρνητικού είναι ο αντίθετός του: $|-8{,}4|=8{,}4=-(-8{,}4)$.
\item $|0|=0$. Η απόλυτη τιμή δεν είναι ποτέ αρνητική.
\end{itemize}

\begin{orismos}{αντίθετοι αριθμοί}
\textbf{Αντίθετοι} είναι δύο ετερόσημοι αριθμοί που έχουν την ίδια απόλυτη τιμή. Ο αντίθετος του $x$ είναι ο $-x$. Για το $0$ ο αντίθετος είναι το $0$. Ισχύει $-(-x)=x$.
\end{orismos}

Δύο σημεία σε ίση απόσταση από το $O$, δεξιά και αριστερά, έχουν αντίθετες τετμημένες: το $+3$ και το $-3$ απέχουν από το $0$ απόσταση $3$ και $|+3|=|-3|=3$.

\tmimasmall{Σύγκριση ρητών}

\begin{kanonas}{σύγκριση ρητών}
Από δύο ρητούς μεγαλύτερος είναι αυτός που βρίσκεται \textbf{δεξιότερα} στον άξονα. Έτσι:
\begin{itemize}[leftmargin=1.2em,itemsep=0.3ex,topsep=2pt]
\item Κάθε θετικός είναι μεγαλύτερος από κάθε αρνητικό. Το $0$ είναι μικρότερο από κάθε θετικό και μεγαλύτερο από κάθε αρνητικό.
\item Από δύο \textbf{θετικούς} μεγαλύτερος είναι αυτός με τη μεγαλύτερη απόλυτη τιμή.
\item Από δύο \textbf{αρνητικούς} μεγαλύτερος είναι αυτός με τη μικρότερη απόλυτη τιμή: $-6{,}8<-3{,}7$, γιατί $|-6{,}8|=6{,}8>3{,}7=|-3{,}7|$.
\end{itemize}
\end{kanonas}

Έτσι $-4{,}75<-3{,}5<-\frac{3}{4}<-0{,}6<0<0{,}6<\frac{3}{4}<3{,}5<4{,}75$.

\begin{prosoxi}
\begin{itemize}[leftmargin=1.2em,itemsep=0.3ex,topsep=2pt]
\item Στους αρνητικούς το «μεγάλο νούμερο» σημαίνει «μικρός αριθμός»: $-10<-3$ παρά το $10>3$.
\item Η απόλυτη τιμή δεν «σβήνει το μείον» από μόνο του: $|-x|=x$ μόνο αν το $x$ είναι θετικό. Η $|-(-3)|=|3|=3$.
\item Η $|x|=3$ έχει δύο λύσεις: $x=3$ και $x=-3$. Η $|x|=-1$ δεν έχει καμία, γιατί η απόλυτη τιμή δεν είναι αρνητική.
\item Ο αντίθετος του $-5$ είναι το $+5$, όχι το $-5$.
\end{itemize}
\end{prosoxi}

\tmima{Λυμένα παραδείγματα}

\paradeigma{Ποια πόλη είναι πιο κρύα}
\tixero Ιωάννινα $-3$°, Φλώρινα $-10$°.\par
\tithelo Ποια θερμοκρασία είναι μικρότερη.\par
\possindeo Συγκρίνω στον άξονα: το $-10$ είναι αριστερότερα από το $-3$. Ή με απόλυτες τιμές: και οι δύο αρνητικοί και $|-10|=10>3=|-3|$, άρα $-10<-3$.\par
\apantisi $-10<-3$: η Φλώρινα είναι πιο κρύα.\par
\elegxos Από το $-10$ μέχρι το $0$ έχουμε $10$ βαθμούς, από το $-3$ μόνο $3$. Συμφωνεί.

\paradeigma{Τοποθετώ και διατάσσω}
\tixero Οι αριθμοί $-2{,}5$, $-0{,}75$, $\frac{1}{2}$, $-\frac{7}{4}$, $3{,}2$.\par
\tithelo Να τους βάλω σε αύξουσα σειρά.\par
\possindeo Γράφω τους δεκαδικούς: $-2{,}5$, $-0{,}75$, $0{,}5$, $-1{,}75$, $3{,}2$. Οι αρνητικοί είναι αριστερά του $0$: από αυτούς μικρότερος είναι αυτός με τη μεγαλύτερη απόλυτη τιμή: $-2{,}5<-1{,}75<-0{,}75$. Μετά οι θετικοί $0{,}5<3{,}2$.\par
\apantisi $-2{,}5<-\frac{7}{4}<-0{,}75<\frac{1}{2}<3{,}2$.\par
\elegxos Στον άξονα, με τη σειρά από αριστερά προς τα δεξιά, εμφανίζονται οι ίδιοι.

\begin{center}
\akgr[1.3]{-3}{4}{-2.5/{-2{,}5},-1.75/{-\frac{7}{4}},-0.75/{-0{,}75},0.5/{\frac{1}{2}},3.2/{3{,}2}}
\end{center}

\paradeigma{Αντίθετος}
\tixero Το σημείο $M$ έχει τετμημένη $+4{,}5$ και το $N$ τετμημένη $-\frac{2}{3}$.\par
\tithelo Οι αντίθετοι.\par
\possindeo Ο αντίθετος του $x$ είναι ο $-x$. Του $+4{,}5$ είναι ο $-4{,}5$. Του $-\frac{2}{3}$ είναι ο $-\left(-\frac{2}{3}\right)=+\frac{2}{3}$.\par
\apantisi $-4{,}5$ και $+\frac{2}{3}$.\par
\elegxos Το άθροισμα απόστασης: $|{+4{,}5}|=|{-4{,}5}|$ και $|-\frac{2}{3}|=|+\frac{2}{3}|$.

\paradeigma{Εξισώσεις με απόλυτη τιμή}
\tixero Οι εξισώσεις $|x|=2$, $|x|=0$, $|x|=-1$.\par
\tithelo Να βρω τους $x$.\par
\possindeo Ψάχνω αριθμούς με απόσταση $2$ από το $0$: $+2$ και $-2$. Με απόσταση $0$: μόνο το $0$. Απόσταση $-1$ δεν υπάρχει.\par
\apantisi $x=2$ ή $x=-2$, $x=0$, αδύνατη.\par
\elegxos $|2|=|-2|=2$.

\paradeigma{Πράξεις με απόλυτες τιμές και αντίθετους}
\tixero Οι παραστάσεις $|-6|\cdot|-2|$, $-|-5|$ και $-(-(-2))$.\par
\tithelo Να τις υπολογίσω.\par
\possindeo $|-6|=6$ και $|-2|=2$, άρα $6\cdot2=12$. Το $|-5|=5$ και ο αντίθετός του είναι $-5$. Ο αντίθετος του $-2$ είναι $2$, και ο αντίθετος του $2$ είναι $-2$.\par
\apantisi $12$, $-5$, $-2$.\par
\elegxos Κάθε δεύτερο «μείον» αναιρεί το προηγούμενο: δύο φορές $-$ επιστρέφει το αρχικό.

\tmima{Ασκήσεις}

\zoni{Το κατάλαβα}{απόλυτη τιμή, αντίθετοι, σύγκριση}

\ex Βρες τις απόλυτες τιμές: α) $|+7|$ \quad β) $|-7|$ \quad γ) $|0|$ \quad δ) $|-2{,}5|$ \quad ε) $\left|-\frac{3}{4}\right|$ \quad στ) $|+0{,}01|$.

\ex Βρες τους αντίθετους των: $5$, $-8$, $0$, $2{,}5$, $-\frac{1}{3}$, $+0{,}6$.

\ex Απλοποίησε: α) $-(-4)$ \quad β) $-(+3)$ \quad γ) $-(-(-2))$ \quad δ) $-|-6|$.

\ex Βρες τον $\alpha$: α) $|\alpha|=5$ \quad β) $|\alpha|=0$ \quad γ) $|\alpha|=\frac{1}{2}$ \quad δ) $|\alpha|=-3$.

\ex Σύγκρινε: α) $-3$ και $-8$ \quad β) $+2$ και $-9$ \quad γ) $-0{,}5$ και $-0{,}05$ \quad δ) $-\frac{7}{2}$ και $-3$ \quad ε) $0$ και $-0{,}001$ \quad στ) $-4$ και $4$.

\ex Γράψε σε αύξουσα σειρά: $-3$, $5$, $-7$, $0$, $2{,}5$, $-0{,}5$, $4$.

\ex Ποιος είναι ο μεγαλύτερος από τους $-1{,}2$, $-1{,}02$, $-1{,}22$ και ποιος ο μικρότερος;

\ex Σωστό ή λάθος; α) $|\alpha|\ge0$ για κάθε $\alpha$. \quad β) $|-5|=-5$. \quad γ) Ο αντίθετος του $-4$ είναι $4$. \quad δ) $-9>-2$. \quad ε) Η $|x|=3$ έχει δύο λύσεις. \quad στ) Ο αντίθετος του $0$ είναι το $0$.

\zoni{Το χρησιμοποιώ}{θερμόμετρα, αποστάσεις, υψόμετρα}

\ex \textbf{Οι θερμοκρασίες.} Μέτσοβο $-5$, Ιωάννινα $-3$, Φλώρινα $-10$, Αθήνα $7$, Θεσσαλονίκη $-1$. Γράψε τις πόλεις από την πιο κρύα στην πιο ζεστή.

\ex \textbf{Οι αποστάσεις.} Στον άξονα της οδού το σχολείο είναι στο $0$ (μονάδα $100$ μ.), ο φούρνος στο $+3$ και το βιβλιοπωλείο στο $-2$. Ποιο είναι πιο κοντά στο σχολείο; Πόση είναι η απόσταση των δύο μαγαζιών μεταξύ τους;

\ex \textbf{Υψόμετρα.} Η Νεκρά Θάλασσα είναι στα $-430$ μ., ο Όλυμπος στα $+2\,917$ μ., η επιφάνεια της θάλασσας στο $0$ και η τάφρος των Μαριανών στα $-11\,034$ μ. Βάλε τα σε αύξουσα σειρά.

\ex \textbf{Το ταμείο.} Η Ζωή έχει $-15$ €, ο Ηλίας $-8$ € και η Ρήνη $+12$ € (αρνητικό σημαίνει ότι χρωστάει). Ποιος χρωστάει περισσότερο; Ποιος έχει το μεγαλύτερο ποσό;

\ex \textbf{Αντίθετες θέσεις.} Το σημείο $K$ έχει τετμημένη $-7$. Βρες το σημείο $\Lambda$ με αντίθετη τετμημένη και πες πόσο απέχουν τα $K$ και $\Lambda$.

\zoni{Το σκέφτομαι}{γιατί το μεγάλο νούμερο είναι μικρός αριθμός}

\ex Ο Βασίλης λέει ότι $-10>-3$, γιατί $10>3$. Εξήγησε το λάθος του με τον άξονα και με το θερμόμετρο.

\ex Γιατί ισχύει πάντα $|\alpha|=|-\alpha|$;

\ex Για ποιους αριθμούς ισχύει $|\alpha|=\alpha$ και για ποιους $|\alpha|=-\alpha$;

\ex Βρες δύο διαφορετικούς αριθμούς $x\neq y$ με $|x|=|y|$. Πόσοι ακέραιοι έχουν απόλυτη τιμή μικρότερη από $2$;

\ex Βρες έναν ρητό ανάμεσα στο $-2$ και στο $-1$ και δύο ρητούς ανάμεσα στο $-0{,}3$ και στο $-0{,}2$. Υπάρχουν άλλοι;

\zoni{Για το σπίτι}{ψάξε θερμοκρασίες}

\ex Βρες τις θερμοκρασίες πέντε πόλεων του κόσμου από τον καιρό και βάλε τες σε σειρά από την πιο κρύα στην πιο ζεστή.
